\chapter{Real analysis, inequalities, and integration}
Before manipulating signals, one must specify their domains, continuity,
integrability, and normed spaces. A sampled waveform is a finite vector;
a continuous-time waveform belongs to a function class. The two are
connected through sampling hypotheses, not merely by notation.

\entry{11}{Triangle inequality}
{In a normed space $\|x+y\|\leq\|x\|+\|y\|$.}
{Bound approximation and reconstruction errors.}
\entry{12}{Cauchy--Schwarz inequality}
{In an inner-product space $|\langle x,y\rangle|
\leq\|x\|\|y\|$; for square-integrable random variables,
$|\Cov(X,Y)|^2\leq\Var(X)\Var(Y)$.}
{Bound Pearson correlation, coherence, and inner-product decoders.}
\entry{13}{Holder's inequality}
{For $1/p+1/q=1$ with $p,q\geq1$, and $f\in L^p$,
$g\in L^q$, $\int|fg|\leq\|f\|_p\|g\|_q$.}
{Establish existence of products and moments.}
\entry{14}{Minkowski's inequality}
{For $p\geq1$, $\|f+g\|_p\leq\|f\|_p+\|g\|_p$.}
{Justify $L^p$ norms on signals and residuals.}
\entry{15}{Bolzano--Weierstrass theorem}
{Every bounded sequence in finite-dimensional Euclidean space
has a convergent subsequence.}
{Understand compactness arguments in parameter estimation.}
\entry{16}{Heine--Borel theorem}
{A subset of $\R^n$ is compact if and only if it is closed and bounded.}
{Provide conditions for existence of finite-dimensional constrained optima.}
\entry{17}{Extreme value theorem}
{Continuous functions on nonempty compact sets attain their
minimum and maximum.}
{Show that certain constrained model-fitting objectives have minimizers.}
\entry{18}{Mean value theorem}
{If $f$ is continuous on $[a,b]$ and differentiable on $(a,b)$,
then some $c\in(a,b)$ satisfies
$f(b)-f(a)=f'(c)(b-a)$.}
{Control local approximation and numerical differentiation errors.}
\entry{19}{Taylor's theorem with remainder}
{Under requisite differentiability, a function equals its order-$k$
Taylor polynomial plus a controlled remainder.}
{Linearize nonlinear neural state equations near operating points.}
\entry{20}{Fundamental theorem of calculus}
{For continuous $f$, $F(x)=\int_a^x f(t)\,dt$
satisfies $F'(x)=f(x)$; suitable primitives recover integrals.}
{Relate firing intensities to integrated expected event counts.}
\entry{21}{Dominated convergence theorem}
{If measurable $f_n\to f$ almost everywhere and
$|f_n|\leq g$ for integrable $g$, then
$\int f_n\to\int f$.}
{Justify exchanging controlled limits and expectations.}
\entry{22}{Monotone convergence and Fatou's lemma}
{For nonnegative measurable increasing $f_n\uparrow f$,
$\int f_n\uparrow\int f$; Fatou gives
$\int\liminf f_n\leq\liminf\int f_n$.}
{Handle nonnegative measures and event-count approximations.}
\entry{23}{Radon--Nikodym theorem}
{If $\nu$ is a sigma-finite measure absolutely continuous
relative to a sigma-finite measure $\mu$, there is a
measurable density $d\nu/d\mu$.}
{Define likelihood and probability density correctly.}
\entry{24}{Banach fixed-point theorem}
{A contraction on a nonempty complete metric space has exactly
one fixed point, reached by iteration from any starting point.}
{Justify convergence of certain recursive estimators under contraction.}

\begin{definition}
A sequence $x_n$ converges to $x$ in a metric space if
for every $\epsilon>0$ there is an $N$ such that
$n\geq N$ implies $d(x_n,x)<\epsilon$.
Convergence in probability, almost surely, and in $L^2$ have
different definitions; none should silently replace another.
\end{definition}
\begin{caution}
The derivative of an estimated signal can amplify measurement noise.
A valid derivative theorem for a smooth function does not prove that a
finite-difference estimate of noisy LFP samples is accurate.
\end{caution}
